\chapter{Pr\'etraitement de texte}
\label{ch:text-preprocessing}

\begin{quote}
\emph{<< Le langage naturel, c'est la donn\'ee la moins structur\'ee qui soit. >>}
\end{quote}

% --- hook ---
<<~Le médecin n'a pas confirmé le diagnostic.~>> Pour un humain, c'est limpide. Pour un modèle de TAL en 2010, c'est un sac de mots sans la moindre relation entre <<~ne~>> et <<~pas~>>. La phrase pouvait aussi bien dire le contraire. Une décennie plus tard, les Transformers ont changé la donne --- mais seulement après que le texte ait été nettoyé, normalisé, tokenisé, vectorisé, dans cet ordre, et avec attention. La qualité de chaque étape se répercute sur tout ce qui suit. Ce chapitre couvre le pipeline standard, avec une attention particulière aux pièges multilingues : beaucoup de cours de TAL supposent implicitement de l'anglais, et écraseront silencieusement les accents d'un texte en français ou les caractères vocaliques d'un texte en arabe.
\section{Pourquoi le texte demande du pr\'etraitement}

Les donn\'ees textuelles ne peuvent pas \^etre fournies directement \`a des mod\`eles num\'eriques. Il faut nettoyer, normaliser, tokeniser, et convertir en repr\'esentations num\'eriques. La qualit\'e du pr\'etraitement d\'etermine directement la qualit\'e de toute t\^ache TAL en aval.

\begin{datatip}
Le pr\'etraitement de texte d\'epend de la langue. R\`egles de tokenisation, stopwords et algorithmes de stemming diff\`erent entre fran\c{c}ais et anglais. Ce chapitre couvre les deux quand c'est pertinent.
\end{datatip}

% ============================================================
\section{Charger des datasets textuels}

\begin{minted}{python}
from sklearn.datasets import fetch_20newsgroups

# 20 Newsgroups : dataset classique de classification de texte
newsgroups = fetch_20newsgroups(subset="train",
                                remove=("headers", "footers", "quotes"))

print(f"Number of documents: {len(newsgroups.data)}")
print(f"Number of categories: {len(newsgroups.target_names)}")
print(f"\nCategories: {newsgroups.target_names}")
print(f"\nSample document:\n{newsgroups.data[0][:300]}")
\end{minted}

\begin{minted}{python}
import pandas as pd

# Cr\'eer un DataFrame pour faciliter la manipulation
df = pd.DataFrame({
    "text": newsgroups.data,
    "category": [newsgroups.target_names[i] for i in newsgroups.target]
})
print(df["category"].value_counts().head(10))
\end{minted}

% ============================================================
\section{Nettoyage de texte}

\begin{minted}{python}
import re

def clean_text(text):
    """Pipeline basique de nettoyage."""
    # Minuscules
    text = text.lower()
    # Supprimer les adresses email
    text = re.sub(r'\S+@\S+', '', text)
    # Supprimer les URL
    text = re.sub(r'http\S+|www\.\S+', '', text)
    # Supprimer les nombres (optionnel : parfois ils comptent)
    text = re.sub(r'\d+', '', text)
    # Supprimer caract\`eres sp\'eciaux, garder lettres et espaces
    text = re.sub(r'[^a-z\s]', '', text)
    # Supprimer les espaces multiples
    text = re.sub(r'\s+', ' ', text).strip()
    return text

# Appliquer au dataset
df["text_clean"] = df["text"].apply(clean_text)
print(df[["text", "text_clean"]].iloc[0])
\end{minted}

\begin{preprocesstip}
Le nettoyage de texte d\'epend de l'ordre. Mettez toujours en minuscules \emph{avant} de supprimer des patterns. Supprimez toujours URL et e-mails \emph{avant} les caract\`eres sp\'eciaux (sinon on d\'etruit les motifs \`a rep\'erer).
\end{preprocesstip}

% ============================================================
\section{Tokenisation}

\begin{minted}{python}
import nltk
nltk.download("punkt_tab", quiet=True)
from nltk.tokenize import word_tokenize, sent_tokenize

text = "Dr. Smith's analysis shows that 42% of patients improved. This is significant!"

# Tokenisation par phrase
sentences = sent_tokenize(text)
print(f"Sentences: {sentences}")

# Tokenisation par mot
tokens = word_tokenize(text)
print(f"Tokens: {tokens}")
print(f"Number of tokens: {len(tokens)}")
\end{minted}

\begin{minted}{python}
# Tokenisation par espaces vs NLTK
text_fr = "L'analyse du Dr. N'Guessan montre que c'est significatif."

# Split par espace : rate les contractions
simple_tokens = text_fr.split()
print(f"Simple: {simple_tokens}")

# NLTK : g\`ere contractions et abr\'eviations
nltk_tokens = word_tokenize(text_fr, language="french")
print(f"NLTK:   {nltk_tokens}")
\end{minted}

% ============================================================
\section{Suppression des stopwords}

\begin{minted}{python}
nltk.download("stopwords", quiet=True)
from nltk.corpus import stopwords

# Stopwords anglais
stop_en = set(stopwords.words("english"))
print(f"English stopwords ({len(stop_en)}): "
      f"{sorted(list(stop_en))[:15]}...")

# Stopwords fran\c{c}ais
stop_fr = set(stopwords.words("french"))
print(f"French stopwords ({len(stop_fr)}): "
      f"{sorted(list(stop_fr))[:15]}...")

# Supprimer les stopwords des tokens
tokens = ["the", "patient", "showed", "significant", "improvement",
          "in", "blood", "pressure", "after", "the", "treatment"]
filtered = [t for t in tokens if t not in stop_en]
print(f"Before: {tokens}")
print(f"After:  {filtered}")
\end{minted}

\begin{warningbox}
Supprimer les stopwords n'est pas toujours b\'en\'efique. En analyse de sentiment, << not good >> devient << good >> apr\`es suppression de << not >>. En topic modeling, c'est g\'en\'eralement utile. \'Evaluez toujours l'impact sur votre t\^ache sp\'ecifique.
\end{warningbox}

% ============================================================
\section{Stemming et lemmatisation}

\begin{minted}{python}
from nltk.stem import PorterStemmer, SnowballStemmer
nltk.download("wordnet", quiet=True)
from nltk.stem import WordNetLemmatizer

stemmer = PorterStemmer()
lemmatizer = WordNetLemmatizer()

words = ["running", "runs", "ran", "studies", "studying",
         "better", "preprocessing", "processed"]

print(f"{'Word':<15} {'Stem':<15} {'Lemma':<15}")
print("-" * 45)
for w in words:
    print(f"{w:<15} {stemmer.stem(w):<15} {lemmatizer.lemmatize(w, 'v'):<15}")
\end{minted}

\begin{minted}{python}
# Stemming fran\c{c}ais avec Snowball
stemmer_fr = SnowballStemmer("french")
mots = ["traitements", "traiter", "traitées", "amélioration",
        "améliorer", "analyse", "analyses", "analyser"]

print("\nFrench stemming:")
for m in mots:
    print(f"  {m:<20} -> {stemmer_fr.stem(m)}")
\end{minted}

\begin{preprocesstip}
Le stemming est rapide mais agressif : << university >> et << universe >> peuvent partager la m\^eme racine. La lemmatisation est plus lente mais linguistiquement correcte : elle renvoie de vrais mots du dictionnaire. Pour la plupart des pipelines NLP, la lemmatisation est pr\'ef\'er\'ee.
\end{preprocesstip}

% ============================================================
\section{Pipeline complet de pr\'etraitement de texte}

\begin{minted}{python}
def preprocess_text(text, language="english"):
    """Pipeline complet de pr\'etraitement de texte."""
    # 1. Nettoyer
    text = text.lower()
    text = re.sub(r'\S+@\S+', '', text)
    text = re.sub(r'http\S+|www\.\S+', '', text)
    text = re.sub(r'[^a-z\sà-ÿ]', '', text)  # garder accentu\'es
    text = re.sub(r'\s+', ' ', text).strip()

    # 2. Tokeniser
    tokens = word_tokenize(text, language=language)

    # 3. Supprimer les stopwords
    stop = set(stopwords.words(language))
    tokens = [t for t in tokens if t not in stop and len(t) > 2]

    # 4. Lemmatiser
    lemmatizer = WordNetLemmatizer()
    tokens = [lemmatizer.lemmatize(t) for t in tokens]

    return tokens

# Appliquer au dataset complet
df["tokens"] = df["text"].apply(preprocess_text)
df["n_tokens"] = df["tokens"].apply(len)
print(df[["category", "n_tokens"]].groupby("category")["n_tokens"].mean()
        .sort_values(ascending=False).head(10))
\end{minted}

% ============================================================
\section{Vectorisation TF-IDF}

\begin{minted}{python}
from sklearn.feature_extraction.text import TfidfVectorizer

# TF-IDF : Term Frequency - Inverse Document Frequency
tfidf = TfidfVectorizer(max_features=5000,
                         stop_words="english",
                         min_df=5, max_df=0.7,
                         ngram_range=(1, 2))

X_tfidf = tfidf.fit_transform(df["text_clean"])
print(f"TF-IDF matrix shape: {X_tfidf.shape}")
print(f"Vocabulary size: {len(tfidf.vocabulary_)}")

# Top termes par score TF-IDF pour le premier document
feature_names = tfidf.get_feature_names_out()
doc_tfidf = X_tfidf[0].toarray().flatten()
top_idx = doc_tfidf.argsort()[-10:][::-1]
print("\nTop TF-IDF terms for document 0:")
for idx in top_idx:
    print(f"  {feature_names[idx]}: {doc_tfidf[idx]:.4f}")
\end{minted}

\begin{minted}{python}
# Bag of Words (alternative plus simple)
from sklearn.feature_extraction.text import CountVectorizer

bow = CountVectorizer(max_features=3000, stop_words="english")
X_bow = bow.fit_transform(df["text_clean"])
print(f"BoW matrix shape: {X_bow.shape}")
\end{minted}

% ============================================================
\section{Expressions r\'eguli\`eres pour l'extraction}

\begin{minted}{python}
import re

# Extraire des donn\'ees structur\'ees d'un texte non structur\'e
texts = [
    "Patient BP: 120/80 mmHg, HR: 72 bpm",
    "Lab result: glucose 142 mg/dL, HbA1c 7.2%",
    "Contact: dr.smith@hospital.org, Tel: +1-555-0123",
]

# Extraire la tension art\'erielle
bp_pattern = r'(\d{2,3})/(\d{2,3})\s*mmHg'
for t in texts:
    match = re.search(bp_pattern, t)
    if match:
        print(f"Systolic: {match.group(1)}, Diastolic: {match.group(2)}")

# Extraire tous les nombres avec unit\'es
num_unit = r'(\d+\.?\d*)\s*(mg/dL|mmHg|bpm|%)'
for t in texts:
    matches = re.findall(num_unit, t)
    for value, unit in matches:
        print(f"  {value} {unit}")
\end{minted}

% ============================================================
\section{Exercices}

\begin{exercise}
\begin{enumerate}
  \item Chargez 20 Newsgroups. Construisez un pipeline complet (nettoyer, tokeniser, retirer stopwords, lemmatiser). Calculez la taille du vocabulaire avant et apr\`es.
  \item Cr\'eez une matrice TF-IDF sur 20 Newsgroups. Trouvez les 10 meilleurs termes pour << sci.space >> et << rec.sport.baseball >>. Sont-ils significatifs ?
  \item \'Ecrivez un pattern regex qui extrait toutes les adresses e-mail d'un texte. Testez-le sur 20 Newsgroups brut.
  \item Comparez le stemming Snowball anglais et fran\c{c}ais sur un document bilingue. Montrez des cas o\`u le stemming produit des racines incorrectes.
  \item Construisez un pipeline de pr\'etraitement pour le fran\c{c}ais. Utilisez les stopwords fran\c{c}ais et le stemmer Snowball fran\c{c}ais. Testez sur 5 phrases sur la science des donn\'ees.
\end{enumerate}
\end{exercise}

% ============================================================
\section{R\'esum\'e du chapitre}

\begin{itemize}
  \item Le pr\'etraitement de texte convertit du texte non structur\'e en variables num\'eriques pour le ML.
  \item Pipeline standard : nettoyer (minuscules, supprimer bruit) $\to$ tokeniser $\to$ retirer stopwords $\to$ stem/lemmatiser $\to$ vectoriser.
  \item Retirer les stopwords aide le topic modeling mais peut nuire \`a l'analyse de sentiment.
  \item La lemmatisation produit de vrais mots et est g\'en\'eralement pr\'ef\'er\'ee au stemming.
  \item TF-IDF pond\`ere les termes par leur importance (fr\'equents dans un document, rares dans le corpus).
  \item Les regex extraient des donn\'ees structur\'ees (nombres, e-mails, motifs) d'un texte non structur\'e.
  \item Le pr\'etraitement multilingue demande tokenizers, listes de stopwords et stemmers sp\'ecifiques \`a la langue.
\end{itemize}
